% --- Archon physics formalization source begin ---
% archon:physics
% archon:covers IPhO2026Problems/problem_IPhO_2026_1_C_2.lean
% archon:source-report reports/ipho_2026/problem_IPhO_2026_1_C_2.source.json
% archon:problem-id IPhO_2026_1
% archon:part-id C.2
% archon:previous-part IPhO_2026_1_C_1
% archon:previous-part-policy natural_language_prerequisite_only

\chapter{Physics problem IPhO\_2026\_1\_C\_2}
\label{ch:IPhO2026Problems_problem_IPhO_2026_1_C_2}

\paragraph{Problem source.}
A photon of angular frequency omega is absorbed by an ozone molecule O3 at rest,
dissociating it into O2 and O.  Let U\_i and U\_f be the ground-state energies of
O3 and O2 and define Delta U = U\_f - U\_i.  The outgoing O2 momentum makes angle
theta with the incident photon.  Treat the oxygen fragments classically and
non-relativistically, take the mass of an oxygen atom to be m, and use photon
momentum p\_gamma = E\_gamma/c = hbar*omega/c.

Current subquestion:
For theta = pi/6, Delta U = 1.10 eV, and m = 16.0 amu, calculate hbar*omega\_min - Delta U in eV.

\paragraph{Current subquestion.}
For theta = pi/6, Delta U = 1.10 eV, and m = 16.0 amu, calculate hbar*omega\_min - Delta U in eV.

\paragraph{Recorded answer/context.}
hbar*omega\_min - Delta U = 2.03e-11 eV.

\paragraph{Figure/image path.}
/root/proposal\_for\_physic/science-mango/ipho\_2026\_source/image/T1\_page-3.png

\paragraph{Reusable previous-part conclusions.}
\begin{itemize}
\item Source C.1, corrected by conservation-law consistency: for
  \(A=2\sin^2\theta+1\) and \(\theta\leq\pi/2\),
  \(\omega_{\min}=3mc^2(1-\sqrt{1-2A\Delta U/(3mc^2)})/(\hbar A)\).
  Policy: natural\_language\_prerequisite\_only; do\_not\_import\_Lean\_output.
\end{itemize}

\paragraph{Formalization target.}
Redraft the existing scaffold so its quoted C.1 prerequisite uses the corrected
factor \(2\).  Preserve the dimensioned constants and momenta, the threshold
meaning, the specified angle, energy and mass data, and the explicit rounding
tolerance.  Leave proof bodies as `by sorry`.

\begin{definition}[DimMass]
  \label{decl:physics:IPhO_2026_1_C_2:DimMass}
  \lean{IPhO2026Problems.IPhO2026_1_C_2.DimMass}
  A mass with no preferred system of units.
\end{definition}
\begin{proof}
  This is the defining declaration; unfolding it gives the stated typed data or relation.
\end{proof}

\begin{definition}[AngularFrequency]
  \label{decl:physics:IPhO_2026_1_C_2:AngularFrequency}
  \lean{IPhO2026Problems.IPhO2026_1_C_2.AngularFrequency}
  An angular frequency, with physical dimension “T⁻¹”, and no preferred system of units.
\end{definition}
\begin{proof}
  This is the defining declaration; unfolding it gives the stated typed data or relation.
\end{proof}

\begin{definition}[DimAction]
  \label{decl:physics:IPhO_2026_1_C_2:DimAction}
  \lean{IPhO2026Problems.IPhO2026_1_C_2.DimAction}
  An action, with physical dimension “M L² T⁻¹”, and no preferred system of units.
\end{definition}
\begin{proof}
  This is the defining declaration; unfolding it gives the stated typed data or relation.
\end{proof}

\begin{definition}[atomicMassUnit]
  \label{decl:physics:IPhO_2026_1_C_2:atomicMassUnit}
  \lean{IPhO2026Problems.IPhO2026_1_C_2.atomicMassUnit}
  \uses{decl:physics:IPhO_2026_1_C_2:DimMass}
  One unified atomic mass unit, represented as a dimensionful mass from its SI readout.
\end{definition}
\begin{proof}
  This is the defining declaration; unfolding it gives the stated typed data or relation.
\end{proof}

\begin{definition}[reducedPlanckConstant]
  \label{decl:physics:IPhO_2026_1_C_2:reducedPlanckConstant}
  \lean{IPhO2026Problems.IPhO2026_1_C_2.reducedPlanckConstant}
  \uses{decl:physics:IPhO_2026_1_C_2:DimAction}
  The reduced Planck constant, represented as a dimensionful action from Physlib's SI value.
\end{definition}
\begin{proof}
  This is the defining declaration; unfolding it gives the stated typed data or relation.
\end{proof}

\begin{definition}[scalarInUnits]
  \label{decl:physics:IPhO_2026_1_C_2:scalarInUnits}
  \lean{IPhO2026Problems.IPhO2026_1_C_2.scalarInUnits}
  The scalar readout of a dimensionful real quantity in a specified system of units.
\end{definition}
\begin{proof}
  This is the defining declaration; unfolding it gives the stated typed data or relation.
\end{proof}

\begin{definition}[siScalar]
  \label{decl:physics:IPhO_2026_1_C_2:siScalar}
  \lean{IPhO2026Problems.IPhO2026_1_C_2.siScalar}
  \uses{decl:physics:IPhO_2026_1_C_2:scalarInUnits}
  The scalar SI readout of a dimensionful real quantity.
\end{definition}
\begin{proof}
  This is the defining declaration; unfolding it gives the stated typed data or relation.
\end{proof}

\begin{definition}[momentumSquaredNorm]
  \label{decl:physics:IPhO_2026_1_C_2:momentumSquaredNorm}
  \lean{IPhO2026Problems.IPhO2026_1_C_2.momentumSquaredNorm}
  The squared Euclidean norm of a spatial momentum in its common working units.
\end{definition}
\begin{proof}
  This is the defining declaration; unfolding it gives the stated typed data or relation.
\end{proof}

\begin{definition}[momentumDot]
  \label{decl:physics:IPhO_2026_1_C_2:momentumDot}
  \lean{IPhO2026Problems.IPhO2026_1_C_2.momentumDot}
  The Euclidean dot product of two momenta expressed in the same working units.
\end{definition}
\begin{proof}
  This is the defining declaration; unfolding it gives the stated typed data or relation.
\end{proof}

\begin{definition}[MakesAngle]
  \label{decl:physics:IPhO_2026_1_C_2:MakesAngle}
  \lean{IPhO2026Problems.IPhO2026_1_C_2.MakesAngle}
  \uses{decl:physics:IPhO_2026_1_C_2:momentumSquaredNorm, decl:physics:IPhO_2026_1_C_2:momentumDot}
  “p” makes the dimensionless angle “θ” with the nonzero reference momentum “q”.
\end{definition}
\begin{proof}
  This is the defining declaration; unfolding it gives the stated typed data or relation.
\end{proof}

\begin{definition}[OzonePhotodissociationSetup]
  \label{decl:physics:IPhO_2026_1_C_2:OzonePhotodissociationSetup}
  \lean{IPhO2026Problems.IPhO2026_1_C_2.OzonePhotodissociationSetup}
  \uses{decl:physics:IPhO_2026_1_C_2:DimMass, decl:physics:IPhO_2026_1_C_2:AngularFrequency}
  The named quantities in Figure 1c.
\end{definition}
\begin{proof}
  This is the defining declaration; unfolding it gives the stated typed data or relation.
\end{proof}

\begin{definition}[DissociationAt]
  \label{decl:physics:IPhO_2026_1_C_2:DissociationAt}
  \lean{IPhO2026Problems.IPhO2026_1_C_2.DissociationAt}
  \uses{decl:physics:IPhO_2026_1_C_2:reducedPlanckConstant, decl:physics:IPhO_2026_1_C_2:scalarInUnits, decl:physics:IPhO_2026_1_C_2:momentumSquaredNorm, decl:physics:IPhO_2026_1_C_2:MakesAngle, decl:physics:IPhO_2026_1_C_2:OzonePhotodissociationSetup}
  A dissociation is kinematically feasible at the given scalar angular-frequency readout and angle.
\end{definition}
\begin{proof}
  This is the defining declaration; unfolding it gives the stated typed data or relation.
\end{proof}

\begin{definition}[ValidOzonePhotodissociationPhysics]
  \label{decl:physics:IPhO_2026_1_C_2:ValidOzonePhotodissociationPhysics}
  \lean{IPhO2026Problems.IPhO2026_1_C_2.ValidOzonePhotodissociationPhysics}
  \uses{decl:physics:IPhO_2026_1_C_2:reducedPlanckConstant, decl:physics:IPhO_2026_1_C_2:scalarInUnits, decl:physics:IPhO_2026_1_C_2:momentumSquaredNorm, decl:physics:IPhO_2026_1_C_2:MakesAngle, decl:physics:IPhO_2026_1_C_2:OzonePhotodissociationSetup, decl:physics:IPhO_2026_1_C_2:DissociationAt}
  The governing physics used for Figure 1c.
\end{definition}
\begin{proof}
  This is the defining declaration; unfolding it gives the stated typed data or relation.
\end{proof}

\begin{definition}[quotedC1ThresholdExpression]
  \label{decl:physics:IPhO_2026_1_C_2:quotedC1ThresholdExpression}
  \lean{IPhO2026Problems.IPhO2026_1_C_2.quotedC1ThresholdExpression}
  \uses{decl:physics:IPhO_2026_1_C_2:reducedPlanckConstant, decl:physics:IPhO_2026_1_C_2:siScalar, decl:physics:IPhO_2026_1_C_2:OzonePhotodissociationSetup}
  The corrected C.1 threshold expression, including the factor \(2\) required by conservation of energy and momentum.
\end{definition}
\begin{proof}
  This is the defining declaration; unfolding it gives the stated typed data or relation.
\end{proof}

\begin{definition}[QuotedPreviousPartC1Result]
  \label{decl:physics:IPhO_2026_1_C_2:QuotedPreviousPartC1Result}
  \lean{IPhO2026Problems.IPhO2026_1_C_2.QuotedPreviousPartC1Result}
  \uses{decl:physics:IPhO_2026_1_C_2:siScalar, decl:physics:IPhO_2026_1_C_2:OzonePhotodissociationSetup, decl:physics:IPhO_2026_1_C_2:quotedC1ThresholdExpression}
  The reusable corrected C.1 conclusion for the forward and backward angle ranges.
\end{definition}
\begin{proof}
  This is the defining declaration; unfolding it gives the stated typed data or relation.
\end{proof}

\begin{definition}[C2NumericalInputs]
  \label{decl:physics:IPhO_2026_1_C_2:C2NumericalInputs}
  \lean{IPhO2026Problems.IPhO2026_1_C_2.C2NumericalInputs}
  \uses{decl:physics:IPhO_2026_1_C_2:atomicMassUnit, decl:physics:IPhO_2026_1_C_2:OzonePhotodissociationSetup}
  The three scalar inputs specified in subquestion C.2, attached to their physical units.
\end{definition}
\begin{proof}
  This is the defining declaration; unfolding it gives the stated typed data or relation.
\end{proof}

\begin{definition}[RoundsTo]
  \label{decl:physics:IPhO_2026_1_C_2:RoundsTo}
  \lean{IPhO2026Problems.IPhO2026_1_C_2.RoundsTo}
  “x” is reported as “reported” to a precision whose half-width is “tolerance”.
\end{definition}
\begin{proof}
  This is the defining declaration; unfolding it gives the stated typed data or relation.
\end{proof}

\begin{theorem}[Physics formalization target]
\label{thm:physics:IPhO_2026_1_C_2:target}
\lean{IPhO2026Problems.IPhO2026_1_C_2.problem_IPhO_2026_1_C_2}
\uses{decl:physics:IPhO_2026_1_C_2:DimMass, decl:physics:IPhO_2026_1_C_2:AngularFrequency, decl:physics:IPhO_2026_1_C_2:DimAction, decl:physics:IPhO_2026_1_C_2:atomicMassUnit, decl:physics:IPhO_2026_1_C_2:reducedPlanckConstant, decl:physics:IPhO_2026_1_C_2:scalarInUnits, decl:physics:IPhO_2026_1_C_2:siScalar, decl:physics:IPhO_2026_1_C_2:momentumSquaredNorm, decl:physics:IPhO_2026_1_C_2:momentumDot, decl:physics:IPhO_2026_1_C_2:MakesAngle, decl:physics:IPhO_2026_1_C_2:OzonePhotodissociationSetup, decl:physics:IPhO_2026_1_C_2:DissociationAt, decl:physics:IPhO_2026_1_C_2:ValidOzonePhotodissociationPhysics, decl:physics:IPhO_2026_1_C_2:quotedC1ThresholdExpression, decl:physics:IPhO_2026_1_C_2:QuotedPreviousPartC1Result, decl:physics:IPhO_2026_1_C_2:C2NumericalInputs, decl:physics:IPhO_2026_1_C_2:RoundsTo}
At \(\theta=\pi/6\), \(\Delta U=1.10\) eV and \(m=16.0\) amu, the corrected
threshold formula makes \(\hbar\omega_{\min}-\Delta U\) round to
\(2.03\times10^{-11}\) eV with the stated tolerance.
\end{theorem}
\begin{proof}
At \(\theta=\pi/6\), the angular factor is \(A=3/2\).  Substitute the supplied
mass and energy into the corrected C.1 expression, convert the dimensioned
quantities to a common SI readout, subtract \(\Delta U\), and divide by one
electronvolt.  The value is approximately \(2.0296693\times10^{-11}\) eV, which
lies inside the requested rounding interval.
\end{proof}
% --- Archon physics formalization source end ---
